\section{Direct bounds for the retarded force}\label{sec:direct}

We first estimate the retarded force by absolute values.  These estimates
control every term except the narrow-angle part of $S_g$ at the scale
needed for the signed bootstrap.  They also give an absolute bound for
the full force, with an additional factor $\sqrt w$.

Retain the decomposition and occupation estimates of the preceding
section.  Let $J\subset[0,\tau]$ have length at most $I$, where
\begin{equation}\label{dir:time-range}
 0<I\le (w/P)^2,
\end{equation}
and let $\mathcal T\subset J$ be any measurable set of times at which
$w/8\le q_X\le8w$.  Every integral below is restricted to
$t\in\mathcal T$.  Estimates for a nonnegative integrand remain valid
if it is multiplied by any measurable weight between zero and one.

\begin{proposition}[Integrated force bounds]\label{dir:direct}
Assume the bootstrap \eqref{eq:bootstrap}.  The bounded data term,
$T_r$, $T_\perp$, $S_b$, and the remaining-sector part of $S_g$ have
combined integral at most
\begin{equation}\label{dir:baseline}
 C\bigl[P\sqrt I+(M+\sqrt A)I S(w)\bigr].
\end{equation}
The integral of the full force satisfies
\begin{equation}\label{dir:absolute}
 \int_{\mathcal T}|K(t)|\,dt
 \le C\sqrt w\bigl[P\sqrt I+(M+\sqrt A)I S(w)\bigr].
\end{equation}
The constant $C$ is independent of $P,w,I,M,A$.
\end{proposition}

We prove the proposition in several steps.  All dyadic sums below have
fixed overlap constants.  The normalization of the radius by
$\max(1,T_0)$ is absorbed into $C$, so that $r\asymp h$.

\subsection{Field budgets and a single bin}

We require field budgets shared by all the momentum and source-scale
bins.  Write $\chi_\phi(t,s,n)$ for the indicator of a fixed enlarged
support of the receiver-angle bin.  The overlap of these supports is
bounded independently of all parameters.  Define
\begin{equation}\label{dir:G-budget}
 \mathcal G_\phi
 =\frac1I\int_{\mathcal T}dt\int_0^t ds\int_{\mathbb S^2}dn\,
       r^2\chi_\phi(t,s,n)G(s,X(t)-rn,n)^2.
\end{equation}
The cone-energy estimate \eqref{eq:cone-energy} gives
$\sum_\phi\mathcal G_\phi\le C$.
In particular this budget does not depend on $p,\theta,\sigma,h$.

For the magnetic field, let $J^*$ be the common enlargement of $J$
from the preceding section, chosen to contain all source times when
$h\le I$; thus $|J^*|\le CI$.  At a point $(s,y)$ with a retarded
receiver time $t_{\rm ret}(s,y)\in(s,\tau]$, put
\[
 n_{\rm ret}(s,y)
 =\frac{X(t_{\rm ret}(s,y))-y}{t_{\rm ret}(s,y)-s}.
\]
Set the indicator in the following definition to zero when the
retarded time does not exist:
\begin{equation}\label{dir:F-budget}
 \mathcal F_\phi
 =\frac1I\int_{J^*}ds\int_{\R^3}dy\,|B(s,y)|^2
    \chi_\phi\bigl(t_{\rm ret}(s,y),s,n_{\rm ret}(s,y)\bigr).
\end{equation}
The retarded time is unique and is a measurable function on its domain.
Bounded overlap and conserved spatial energy \eqref{eq:energy} imply
$\sum_\phi\mathcal F_\phi\le C$.  Again the budget is independent of
$p,\theta,\sigma,h$.  The Jacobian $r^2D$ of the map
$(t,n)\mapsto X(t)-(t-s)n$, together with $D\asymp\phi^2$, gives
\begin{equation}\label{dir:magnetic-budget}
 \int_{\mathrm{bin}} r^2\phi^2|B|^2\,dt\,ds\,dn
 \le C(I+h)\mathcal F,
 \qquad
 \mathcal F=
 \begin{cases}
  \mathcal F_\phi,&h\le I,\\
  1,&h>I.
 \end{cases}
\end{equation}
For $h>I$ the source times lie in an interval of length $O(I+h)$,
so conserved spatial energy proves the second case directly.
The budgets in \eqref{dir:G-budget}--\eqref{dir:F-budget} may always
be enlarged by fixed factors when supports are enlarged.

For a fixed bin, let $N=\int_{\mathrm{bin}}f\,dv$ and denote its hit
measure by $\mathsf M$.  By \eqref{occ:density}, the angular cap bound,
and the definition of $\mathsf M$,
\begin{equation}\label{dir:density-integrals}
 \begin{split}
 N&\le Cp^3\nu^2,\qquad
 \int_{\mathrm{bin}}dt\,ds\,dn\le CIh\phi^2,\\
 \int N\,dt\,ds\,dn&\le \frac{C\mathsf M}{h^2\phi^2},\\
 \int N^2\,dt\,ds\,dn
 &\le C\min\left(Ih\phi^2p^6\nu^4,
              \frac{p^3\nu^2\mathsf M}{h^2\phi^2}\right).
 \end{split}
\end{equation}
Here and below $N$ vanishes outside the indicated time and spatial
restrictions.  In the small sector $\nu=\theta$ and the $\sigma$
partition is omitted.

The force majorants on this bin are, up to fixed constants,
\[
 T_r=\frac{f}{r^2p^2\theta^2},\qquad
 T_\perp=\frac{f\phi}{r^2p^2\theta^3},\qquad
 S_g=\frac{f\phi G}{rp\theta^2},\qquad
 S_b=\frac{f\phi|B|}{rp}.
\]
Using $dy=r^2ds\,dn$ yields the following bounds for their integrals.
Both columns apply on every bin: the near bound uses the pointwise density
estimate, while the far bound also uses the occupation measure. Their names
indicate which estimate is useful at smaller or larger radii.
\begin{equation}\label{dir:near-far}
\renewcommand{\arraystretch}{1.5}
\begin{array}{c|c|c}
 &\text{near bound}&\text{far bound}\\ \hline
 T_r&Ihp(\nu/\theta)^2\phi^2&-\\
 T_\perp&Ihp\theta^{-3}\nu^2\phi^3
         &\mathsf M/(h^2\phi p^2\theta^3)\\
 S_g&I\sqrt h\,p^2(\nu/\theta)^2\phi^2\sqrt{\mathcal G_\phi}
     &\sqrt I\sqrt{p\mathsf M}\,\nu(\theta^2h)^{-1}
        \sqrt{\mathcal G_\phi}\\
 S_b&\sqrt{I(I+h)}\sqrt h\,p^2\nu^2\phi\sqrt{\mathcal F}
     &\sqrt{I+h}\sqrt{p\mathsf M}\,\nu(\phi h)^{-1}
        \sqrt{\mathcal F}.
\end{array}
\end{equation}
The first two rows follow from \eqref{dir:density-integrals}.
For $S_g$, apply Cauchy--Schwarz to $rGN$: the first squared factor
is bounded by $I\mathcal G_\phi$ and the second by either bound for
$\int N^2$ in \eqref{dir:density-integrals}.  The prefactor is
$\phi/(p\theta^2)$.  For $S_b$, use $r\phi|B|$ as the first factor,
apply \eqref{dir:magnetic-budget}, and use the prefactor $1/p$.
This proves every entry of the table.

\subsection{Summing the radius}

Recall the occupation factor
\[
 R=1+\sqrt{I/\lambda}\sqrt{\mathcal E},
 \qquad
 \mathcal E=
 \begin{cases}
  1,&\theta\ge\kappa\phi,\\
  \mathcal E_{p\phi},&\theta<\kappa\phi,
 \end{cases}
 \qquad \sum_{p,\phi}\mathcal E_{p\phi}\le C.
\]
It is independent of $h$, and also of $\theta$ in the small sector.
The refined occupation estimate \eqref{occ:refined} gives
\begin{equation}\label{dir:occupation}
 \mathsf M\le C h k_1R^2/p
\end{equation}
when $h\le I$, except in deficit bins with $h\le\lambda$.
For $h>I$ the same inequality follows from \eqref{occ:energy},
since $R,k_1\ge1$.

For positive $a,b,\alpha,\beta$, the elementary dyadic estimate
\begin{equation}\label{dir:balance}
 \sum_{h\in2^{\mathbb Z}}
       \min(ah^\alpha,bh^{-\beta})
 \le C_{\alpha,\beta}
       a^{\beta/(\alpha+\beta)}b^{\alpha/(\alpha+\beta)}
\end{equation}
follows by splitting at $h=(b/a)^{1/(\alpha+\beta)}$ and summing
geometric series on both sides.  It remains an upper bound when the
allowed radii form only a subset of the dyadic scales.
Substituting \eqref{dir:occupation} into \eqref{dir:near-far} and
applying \eqref{dir:balance} gives
\begin{equation}\label{dir:balanced}
 \begin{aligned}
 T_\perp:&\quad
 C\sqrt I\,R p^{-1}\theta^{-3}\nu\phi\sqrt{k_1},\\
 S_g:&\quad
 C I^{3/4}R^{1/2}p\phi\nu^{3/2}\theta^{-2}
       k_1^{1/4}\sqrt{\mathcal G_\phi},\\
 S_b\ (h\le I):&\quad
 C I^{3/4}R^{1/2}p\nu^{3/2}k_1^{1/4}
       \sqrt{\mathcal F_\phi}.
 \end{aligned}
\end{equation}
The first line balances $h$ with $h^{-1}$; the other two balance
$h^{1/2}$ with $h^{-1/2}$.  The exceptional deficit radii are excluded
from these balances and are treated below.

For $S_b$ with $h>I$, use instead the time-change bound
\eqref{occ:time}, namely $\mathsf M\le CI\phi^2/\theta^2$.
The corresponding minimum in \eqref{dir:near-far} is
\[
 C\sqrt I\min\left(hp^2\nu^2\phi,
                     h^{-1/2}\sqrt p\,\nu/\theta\right).
\]
Its radius sum is at most
\begin{equation}\label{dir:large-radius}
 C\sqrt I\,p\nu^{4/3}\theta^{-2/3}\phi^{1/3}.
\end{equation}
If one of $\theta,\phi$ is much larger than the other, $\sigma$ has
bounded multiplicity or is absent, and the angular factor is at most
$\theta^{2/3}\phi^{1/3}$.  If $\theta\asymp\phi$, it is at most
$C\theta(\sigma/\theta)^{4/3}$, with $\sigma\lesssim\theta$.
The angular sums are geometric in both cases.  Summing $p$ proves
that all large-radius $S_b$ terms cost at most $CP\sqrt I$.

In an exceptional deficit bin,
$h\le\min(I,\lambda)$ and $\lambda\le CP^{-2}$ by
\eqref{occ:lambda}.  Summing the near bounds over these radii gives,
per remaining bin,
\begin{equation}\label{dir:deficit}
 \begin{aligned}
 S_g+S_b&\le CI\frac{p^2}{Pm^2},\\
 T_\perp&\le
 \begin{cases}
  CIp/P^2,&\theta\ge\kappa\phi,\\
  CIp^2/P^2,&\theta<\kappa\phi.
 \end{cases}
 \end{aligned}
\end{equation}
Indeed, in the remaining sector $\phi\lesssim\theta$ and
$\nu\lesssim m^{-1}$, whereas in the small sector
$\nu=\theta$, $\phi\lesssim m^{-1}$, and $\theta^{-1}\le p$.
The individual field budgets are bounded by their total budgets.
There are at most $CL^4$ remaining bins.  Since
\[
 \frac{p^2}{Pm^2}\le \frac1P+\frac{P}{w^2},\qquad
 \frac{p^2}{P^2}\le1,
\]
the ratio of their total cost to $I S(w)$ is at most
\[
 CL^3\left(\frac{w}{P^3}+\frac1{Pw}+\frac{w}{P^2}\right)
 \le \frac{CL^3}{P}.
\]
Thus these terms are absorbed uniformly into $CI S(w)$.

The radius sum for $T_r$ is also immediate from its near bound.
When $\theta\asymp\phi$, summing $(\nu/\theta)^2$ over $\sigma$
is a geometric sum; otherwise $\sigma$ has bounded multiplicity or
is omitted.  Summing $h$, the $O(L)$ values of $\theta$, then
$\phi^2$ and $p$, gives
\begin{equation}\label{dir:radial-transport}
 \int T_r\le CIPL\le CI S(w).
\end{equation}

\subsection{Angular and momentum sums}

The radius sums are complete, including the deficit exception and the
magnetic terms with $h>I$. It remains to sum the angular and momentum
factors in \eqref{dir:balanced}. We keep the small-sector $S_g$
contribution separate because it yields the absolute-force bound with
its additional factor $\sqrt w$. The following elementary momentum
sums will be used repeatedly:
\begin{equation}\label{dir:momentum-sums}
 \begin{aligned}
 \sum_p p&\le CP,&
 \sum_p\frac p m&\le C(L+P/w),&
 \sum_p\frac{\sqrt p}{m}&\le C(1+\sqrt P/w),\\
 \sum_p p\sqrt{P/m}&\le CP\sqrt{P/w},&&
 \sum_p p^{3/2}\sqrt{P/m}&\le CP^2/\sqrt w.
 \end{aligned}
\end{equation}
To verify them, split at $p=w$.  For $p\le w$ one has $m=p$;
for $p>w$ one has $m=w$.  Each resulting positive-power sum is
geometric, while a constant summand costs at most $CL$.

Write $\omega=\sigma$ in the remaining sector and
$\omega=\phi$ in the small sector.  The definition of $\lambda$
gives the exact parameter dependence
\begin{equation}\label{dir:R-expansion}
 \begin{split}
 R-1
 &\le C\sqrt I\frac Pm
       \left(\frac M\omega+
                     \frac{\sqrt{AL}}{\sqrt\omega}\right)
       \sqrt{\mathcal E},\\
 R^{1/2}
 &\le C\left[1+I^{1/4}\sqrt{P/m}
       \left(\frac{\sqrt M}{\sqrt\omega}+
                     \frac{(AL)^{1/4}}{\omega^{1/4}}\right)
       \right].
 \end{split}
\end{equation}
For the second line we used $\mathcal E\le C$; the first line retains
the joint small-sector budget when it is needed.

\paragraph{The small-sector transport term.}
Here $\nu=\theta$ and $k_1\asymp1$.  Summing the first line of
\eqref{dir:balanced} over $\theta\ge1/p$ gives
\begin{equation}\label{dir:small-budget}
 C\sqrt I\,R p\phi.
\end{equation}
The term with $R$ replaced by one sums to $CP\sqrt I$.
The $M$ term in \eqref{dir:R-expansion} costs
\[
 CIM P\sum_{p,\phi}\frac p m\sqrt{\mathcal E_{p\phi}}
 \le CIM P\bigl(L+\sqrt L\,P/w\bigr).
\]
Indeed, Cauchy--Schwarz and the joint budget reduce the sum to the
square root of $\sum_{p,\phi}(p/m)^2$, which is at most
$C[L^2+L(P/w)^2]$ before taking the square root.  For the
$\sqrt{AL}$ term the weight has an additional factor $\sqrt\phi$;
its squared sum is at most $C[L+(P/w)^2]$.  That contribution is
therefore at most
\[
 CI\sqrt{AL}\,P(\sqrt L+P/w)
 =CI\sqrt A\,P(L+\sqrt L\,P/w).
\]
Since $1\le w\le P$ and $L\ge1$, both parameter terms are bounded
by $C(M+\sqrt A)I S(w)$.

\paragraph{The remaining-sector transport term.}
Split into $\theta>16\phi$, where $\sigma\asymp\theta$, and
$\theta\asymp\phi$, where $\sigma\lesssim\theta$.  Fixed comparison
constants may depend on $\kappa$.  The angular factors in the plain
term of the first line of \eqref{dir:balanced} are respectively
\[
 p^{-1}\theta^{-2}\phi,
 \qquad p^{-1}\theta^{-1}(\sigma/\theta)^{1/2}.
\]
Their sums are bounded uniformly for each $p$: sum $\phi\lesssim
\theta$ in the first expression, and $\sigma\lesssim\theta$ in the
second, then use $\theta\ge1/p$.  Thus the plain contribution is at
most $CL\sqrt I\le CP\sqrt I$.

After extracting $CIM P(p/m)$, the $M$ terms have angular factors
\[
 (p\theta)^{-2}\frac\phi\theta,
 \qquad (p\theta)^{-2}(\theta/\sigma)^{1/2}.
\]
The second $\sigma$ sum is at most $C\sqrt{m\theta}\le
C\sqrt{p\theta}$ because $\sigma\gtrsim1/m$; hence both remaining
angular sums are bounded.  After extracting
$CI\sqrt{AL}\,P\sqrt p/m$, the corresponding factors are
\[
 (p\theta)^{-3/2}\frac\phi\theta,
 \qquad (p\theta)^{-3/2}.
\]
The number of allowed $\sigma$ values in the second case is at most
$C[1+\log_+(m\theta)]\le C\sqrt{p\theta}$ on a nonempty bin, so
these angular sums are bounded as well.  Here
$\log_+s=\max(0,\log s)$.  By
\eqref{dir:momentum-sums}, the two parameter contributions are at
most
\[
 CI M P(L+P/w),\qquad
 CI\sqrt{AL}\,P(1+\sqrt P/w).
\]
Both are bounded by $C(M+\sqrt A)I S(w)$.

\paragraph{The remaining-sector good-field term.}
Apart from $CI^{3/4}pR^{1/2}\sqrt{\mathcal G_\phi}$,
the angular factors for $S_g$ in the same two cases are
\[
 \sqrt\theta\,\frac\phi\theta,
 \qquad \sqrt\theta(\sigma/\theta)^{5/4}.
\]
At fixed $\phi$ their $\theta,\sigma$ sums are at most
$C\sqrt\phi$.  Consequently the plain term is at most
$CI^{3/4}P\le CP\sqrt I$, by Cauchy--Schwarz with
$\sum_\phi\mathcal G_\phi\le C$ and $\sum_\phi\phi\le C$.
Dividing the displayed factors by $\sqrt\sigma$ leaves, up to
constants,
\[
 \phi/\theta,
 \qquad (\sigma/\theta)^{3/4},
\]
whose sums at fixed $\phi$ are bounded.  Division by
$\sigma^{1/4}$ gives no larger bound, since $\sigma\le1$.
There are at most $CL$ receiver-angle bins, so
$\sum_\phi\sqrt{\mathcal G_\phi}\le C\sqrt L$.
The parameter terms in \eqref{dir:R-expansion} therefore cost at most
\begin{equation}\label{dir:field-parameters}
 CI\bigl(\sqrt M+(AL)^{1/4}\bigr)
       P\sqrt{P/w}\sqrt L
 \le C(M+\sqrt A)I S(w).
\end{equation}
The last inequality follows from $w\le P$, $L\ge1$, and $M,A\ge1$.

\paragraph{The magnetic remainder for $h\le I$.}
For $S_b$ the base angular factors in the two remaining-sector subcases and
in the small sector are respectively
\[
 \theta^{3/2},\qquad
 \theta^{3/2}(\sigma/\theta)^{5/4},\qquad
 \theta^{3/2}.
\]
Their sums at fixed $\phi$ are bounded.  This remains true after
division by $\sqrt\sigma$ in the remaining-sector subcases, giving
$\theta$ and $\theta(\sigma/\theta)^{3/4}$, or by $\sqrt\phi$ in
the small sector, where the $\theta$ sum is at most $C\phi$.
The quarter-power divisions give no larger bound.  Using
$\sum_\phi\sqrt{\mathcal F_\phi}\le C\sqrt L$ thus gives the
parameter estimate \eqref{dir:field-parameters} again.
The plain term is at most $CI^{3/4}P\sqrt L$.  Uniformly in $I$,
\begin{equation}\label{dir:plain-magnetic}
 I^{3/4}\sqrt L\le \sqrt I+IL,
\end{equation}
as follows from $x\le1+x^2$ with $x=I^{1/4}\sqrt L$.
Since $PL\le S(w)$, this completes the bound for $S_b$.

\paragraph{The small-sector good-field term.}
This term is needed only for the absolute estimate
\eqref{dir:absolute}.  Its base angular factor is
$\phi/\sqrt\theta$, whose $\theta$ sum is at most
$C\sqrt p\,\phi$.  The plain contribution is therefore bounded by
$CI^{3/4}P^{3/2}$.  The time restriction \eqref{dir:time-range} gives
\[
 I^{3/4}P^{3/2}\le \sqrt w\,P\sqrt I.
\]
For the two parameter terms the $\phi$ factors after this sum are
$\phi^{1/2}$ and $\phi^{3/4}$.  Their sums against
$\sqrt{\mathcal G_\phi}$ are bounded by Cauchy--Schwarz and
geometric summation.  The last momentum sum in
\eqref{dir:momentum-sums} gives a total of at most
\[
 CI\bigl(\sqrt M+(AL)^{1/4}\bigr)P^2/\sqrt w
 \le C\sqrt w(M+\sqrt A)I S(w).
\]
The exceptional bins were already included in \eqref{dir:deficit}.

Combining these bounds with \eqref{dir:large-radius},
\eqref{dir:deficit}, and \eqref{dir:radial-transport} proves
\eqref{dir:baseline}.  The data force contributes at most $CI$,
which is bounded by $CP\sqrt I$ since $I\le1$.
Finally the force representation is bounded by these nonnegative
majorants; adding the small-sector $S_g$ estimate and using $w\ge1$
proves \eqref{dir:absolute}.  This proves
Proposition~\ref{dir:direct}.
